\section{Protocol Mismatch Manufactures Evidence: A Retraction}
\label{sec:retraction}

Our first attempt at this comparison produced the opposite result, and the way it failed is the
most transferable finding in this paper.

That attempt computed $\mathrm{score}_5$ with the \textsc{UnseenEqClass} split serving as its own
neighbour pool, and normalised by $k$. Under that procedure the non-learned floor looked dramatic: a
bag or an untrained encoder matched or beat the StructEmb ablation on 13 of 14 corpora, beat the
headline SemEmb system on 3, and reached $69.4$ on \textsc{Bool8} where EqNet reports $58.1$. Had we
stopped there, we would have published a claim that a zero-parameter baseline undercuts a published
leaderboard.

It was wrong, and the direction of the error was not random. Three differences from the published
procedure, neighbour pool, normalisation, and query construction
(Appendix~\ref{app:retraction}), \emph{each} pushed our numbers up. On \textsc{Bool8} the
published protocol searches $257{,}784$ candidate neighbours where ours searched $13{,}859$, and
more distractors can only lower a score. Reproducing all three moves the comparison from ``matches
or beats StructEmb on 13 of 14'' to 5 of 14, and from ``beats SemEmb on 3'' to none.

\paragraph{Why this is not just our bug.} We had done what a careful auditor is supposed to do. We
reported \emph{both} non-learned baselines, under a \emph{single, internally consistent}
protocol, satisfying our own checklist, and still produced a badly misleading comparison,
because the protocol was ours and the numbers beside it were someone else's. Internal consistency
is not the relevant property when a baseline is placed in a table next to published results; what
matters is whether the baseline was computed under \emph{the same evaluation} as the numbers it is
being compared against. Every non-learned floor, hypothesis-only baseline, or majority-class
control reported against an external leaderboard is exposed to this error, and the error has a
direction: the shortcuts that make a protocol convenient to re-implement (evaluate within the
test split, normalise by $k$) are the ones that inflate the floor.

We keep the superseded numbers on record rather than deleting them, and we add a checklist step for
it (step~7, \S\ref{sec:checklist}): verify the evaluation matches what you compare against.
